%Document Class
\documentclass[a4paper,12pt,reqno]{amsart}


  \headheight0.6in
  \headsep22pt
  \textheight23cm
  \topmargin-1.7cm
  \oddsidemargin 0.5cm
  \evensidemargin0.5cm
  \textwidth15.3cm

%Packages 
\usepackage{amscd,amssymb,amsmath,amsfonts,amsthm, mathrsfs,xspace}
\usepackage{graphicx}
\usepackage{verbatim,enumerate,paralist}
\usepackage{textcomp}
%\usepackage[authoryear]{natbib}
\usepackage[pagebackref,colorlinks]{hyperref}
\hypersetup{%
  urlcolor=blue,linkcolor=blue,citecolor=blue
}
\usepackage{pdfpages}
%\usepackage{pdfsync}%  remove this when finalised

% import these fonts by hand here to avoid clashes with usual blackboard bold usage.
%\pdfmapfile{+bbold.map}
\newcommand{\bbefamily}{\fontencoding{U}\fontfamily{bbold}\selectfont}
\newcommand{\textbbe}[1]{{\bbefamily #1}}
\DeclareMathAlphabet{\mathbbe}{U}{bbold}{m}{n}

% Macro to typeset part of a math formula at a bigger size
\def\makebigger#1#2#3{\hbox{#1$\mathsurround=0pt #2{#3}$}}
\def\bigger#1#2{{\relax\mathpalette{\makebigger#1}{#2}}}

% Macro for the category Delta
\newcommand{\DelCat}{{\bigger\large{\mathbbe{\Delta}}}}
\newcommand{\DelCAT}{{\bigger\Huge{\mathbbe{\Delta}}}}


\usepackage[T1]{fontenc}
\usepackage[all,2cell,poly,knot,tips]{xy}
\UseAllTwocells
% \CompileMatrices

\def\blue{\color{blue}}
\def\red{\color{red}}
\def\black{\color{black}}

%\usepackage{xypdf} % may not be necessary
\setlength{\marginparwidth}{2cm}
\newcommand{\missing}[1]{\marginpar{\color{red}missing:\\#1}}

%Theoretic Environment Setup
\newcounter{Definitioncount}
\setcounter{Definitioncount}{0}

\newtheorem*{Theorem}{Theorem}% no counter
\newtheorem{theorem}{Theorem}[section] % with counter
\newtheorem{lemma}[theorem]{Lemma}
\newtheorem{proposition}[theorem]{Proposition}
% \newtheorem*{Proposition}{Proposition}
\newtheorem{corollary}[theorem]{Corollary}
% \newtheorem*{Corollary}{Corollary}% no counter

\newtheorem*{Conjecture}{Conjecture}% no counter
\newtheorem*{Lemma}{Lemma}% no counter

\theoremstyle{definition}
\newtheorem*{Definition}{Definition}% no counter
\newtheorem*{Observation}{Observation}% no counter
\newtheorem*{Observations}{Observations}% no counter
\newtheorem*{Remark}{Remark}% no counter
\newtheorem{remark}[theorem]{Remark}
\newtheorem*{Example}{Example}% no counter
\newtheorem{example}[theorem]{Example}

\renewcommand{\theexample}{\arabic{example}}

\newtheoremstyle{fact}{\bigskipamount}{\medskipamount}{\upshape}{}{\itshape}{. }{ }{Fact}
\theoremstyle{fact}
\newtheorem*{Fact}{Fact}% no counter

\newtheoremstyle{genquest}{\bigskipamount}{\medskipamount}{\upshape}{}{\itshape}{. }{ }{General Question}
\theoremstyle{genquest}
\newtheorem*{GenQuest}{General Question}% no counter

% use \newtheoremstyle to underline the header
\newtheoremstyle{step}{2\bigskipamount}{\medskipamount}{\upshape}{}{\itshape}{. }{ }{\underline{Step~\thestep}}
\theoremstyle{step}
\newtheorem{step}{Step}[subsection]
\renewcommand{\thestep}{\arabic{step}}

\newcommand{\lra}{\longrightarrow}
\newcommand{\lla}{\longleftarrow}
\newcommand{\Lra}{\Longrightarrow}
\newcommand{\Lla}{\Longleftarrow}
\newcommand{\ldual}[1]{\mathord{{\let\nolimits\relax\sideset{^\wedge}{}{#1}}}}
\newcommand{\laction}[2]{\mathord{{\let\nolimits\relax\sideset{^{#1}}{}{#2}}}}
\newcommand{\conj}[2]{\mathord{{\let\nolimits\relax\sideset{^{#1}}{}{#2}}}}
\newcommand{\xylongto}[2][2pt]{{\UseTips{}\xymatrix @C+#1{ \ar[r]^{#2}&}}}
\newcommand{\ola}{\overleftarrow}
\newcommand{\ora}{\overrightarrow}

\DeclareMathOperator{\im}{im}

%% Script Shortcuts 
\def\CA{{\mathscr A}}
\def\CB{{\mathscr B}}
\def\CC{{\mathscr C}}
\def\CD{{\mathscr D}}
\def\CE{{\mathscr E}}
\def\CF{{\mathscr F}}
\def\CG{{\mathscr G}}
\def\CH{{\mathscr H}}
\def\CI{{\mathscr I}}
\def\CJ{{\mathscr J}}
\def\CK{{\mathscr K}}
\def\CL{{\mathscr L}}
\def\CM{{\mathscr M}}
\def\CN{{\mathscr N}}
\def\CO{{\mathscr O}}
\def\CP{{\mathscr P}}
\def\CQ{{\mathscr Q}}
\def\CR{{\mathscr R}}
\def\CS{{\mathscr S}}
\def\CT{{\mathscr T}}
\def\CU{{\mathscr U}}
\def\CV{{\mathscr V}}
\def\CW{{\mathscr W}}
\def\CX{{\mathscr X}}
\def\CY{{\mathscr Y}}
\def\CZ{{\mathscr Z}}
\def\ordm{{\mathbf{m}}}
\def\ordn{{\mathbf{n}}}
\def\ordk{{\mathbf{k}}}


\newcommand{\ox}{\otimes}
\newcommand{\x}{\times}

\newcommand{\op}{\ensuremath{^{\textnormal{op}}}}


\def\theequation{{\thesection.\arabic{equation}}}


\begin{document}

\title{Morita equivalence of partial maps and strong epimorphisms}

\author{Stephen Lack}
\address{Department of Mathematics, Macquarie University, NSW 2109, Australia}
\email{steve.lack@mq.edu.au}

\author{Ross Street}
\address{Department of Mathematics, Macquarie University, NSW 2109, Australia}
\email{ross.street@mq.edu.au}

\thanks{Both authors gratefully acknowledge the support of the Australian Research Council Discovery Grant DP130101969; Lack acknowledges with equal gratitude the support of an Australian Research Council Future Fellowship}

\date{\today}
\maketitle

\tableofcontents

\section{Introduction}

The intention of this paper is to give a rather pretty categorical proof
of an equivalence of categories that seems of interest in 
enumerative combinatorics as per \cite{Species},
representation theory as per \cite{CEF2012}, 
and homotopy theory as per \cite{Berger1996}.

Our results involve a category $\CA$ equipped with a fairly general 
factorization system $(\CE,\CM)$ in the sense of \cite{FreydKelly}.
Pullbacks of $\CM$s along arbitrary morphisms should exist, leading to
a category $\CP$ of partial maps in $\CA$. We obtain a fully faithful
functor $$\widehat{(-)} : [\CE, \mathrm{Mod}^R] \lra [\CP,\mathrm{Mod}^R]$$ 
with a right adjoint and monomorphic counit. 
If $\CA$ is finitely $\CM$-well powered
(that is, if every object only has finitely many $\CM$-subobjects), 
we prove that this functor is an equivalence.

Let $\mathfrak{S}$ be the groupoid of finite sets and bijective functions.
Let $\mathrm{FI}\sharp$ denote the category of finite sets and injective partial functions.
Let $\mathrm{Mod}^R$ denote the category of left modules over the ring $R$.
Our purpose is to understand and generalize the classification theorem for $\mathrm{FI}\sharp$-modules appearing as Theorem 2.24 of \cite{CEF2012}, which provides an equivalence between the category of functors
$\mathfrak{S} \lra \mathrm{Mod}^R$ and the category of functors
$\mathrm{FI}\sharp \lra \mathrm{Mod}^R$. 
Their result gives a new viewpoint on representations of the symmetric groups,
also a new viewpoint on Joyal species \cite{Species, AnalFunct}.

\section{The kernel module}\label{km}

Let us begin with any category $\CA$ equipped with a 
prefactorization system $(\CE,\CM)$ in the sense of \cite{FreydKelly}.
We regard both $\CE$ and $\CM$ as subcategories of $\CA$
having the same objects as $\CA$ but only the respective 
morphisms of their class.  
Recall that, for a morphism $e : B \lra A$ in $\CE$ and a morphism $m : X \lra Y$
in $\CM$, the following square is a pullback.
\begin{eqnarray}\label{strepi}
\begin{aligned}
\xymatrix{
\CA(A,X) \ar[rr]^-{\CA(e,X)} \ar[d]_-{\CA(A,m)} && \CA(B,X) \ar[d]^-{\CA(B,m)} \\
\CA(A,Y) \ar[rr]_-{\CA(e,Y)} && \CA(B,Y)}
\end{aligned}
\end{eqnarray}

An {\em $\CM$-subobject} of an object $X$ of $\CA$ is an isomorphism
class, in the slice category $\CA /X$, of objects $m : U \lra X$ with $m\in \CM$.    

Assume that the pullback of $\CM$-morphisms along arbitrary morphisms exist. 
A span $f = (X\stackrel{f_0}\lla U \stackrel{f_1}\lra Y)$ is called an 
{\em ($\CM$-)partial map} $f : X\lra Y$ when $f_0\in \CM$. 

Let $\CP = \mathrm{Par}_{\CM}\CA$ 
denote the category whose objects are all those of $\CA$ and 
whose morphisms are isomorphism classes $[f]$ of partial maps. 
Composition is that of spans: that is, by pullback.
There are functors $(-)^* : \CM^{\mathrm{op}} \lra \CP$
and $(-)_* : \CA \lra \CP$ which are the identity on objects.
For $m:U\lra X$ in $\CM$, define $m^*= [m,U,1_U] : X\lra U$.
For any $f: X \lra Y$ in $\CA$, define $f_*=[1_X,X,f] : X \lra Y$.
Every partial map $f = (X\stackrel{f_0}\lla U \stackrel{f_1}\lra Y)$ has
a factorization $$[f] = f_{1*}\circ f_0^* \ .$$ 
For every pullback
\begin{equation}\label{pbm}
\xymatrix{
P \ar[rr]^-{g} \ar[d]_-{n} && B \ar[d]^-{m} \\
A \ar[rr]_-{f} && C}
\end{equation}
in $\CA$ with $m$ a monomorphism, the square
\begin{equation}
\xymatrix{
P \ar[rr]^-{g_*}  && B  \\
A \ar[rr]_-{f_*} \ar[u]^-{n^*} && C \ar[u]_-{m^*}}
\end{equation}
commutes in $\CP$. 

There is analogue, for partial maps in place of general spans, of the result in
\cite{Lindner1976} about Mackey functors.
There is an equivalence between functors $T : \CP \lra \CX$
and pairs $(T^*,T_*)$ of functors 
$T^* : \CM^{\mathrm{op}} \lra \CX$
and $T_* : \CA \lra \CX$
which agree on objects and satisfy the B\'enabou-Roubaud-Chevalley-Beck condition on pullbacks of the form \eqref{pbm}. 
Indeed $T^*=T\circ (-)^*$ and $T_*=T\circ (-)_*$. 

Let $\mathrm{Set}$ denote the category of sets and functions. 
We write $\mathrm{Par}\mathrm{Set}$ for $\mathrm{Par}_{\CM}\mathrm{Set}$ 
when $\CM$ consists of the injective functions.

Notice that, if all morphisms in $\CA$ are in $\CM$ (such as in the
category of finite sets and injective functions), 
then $\CE$ consists of the invertible morphisms; it is a groupoid.

There is a ``kernel'' functor 
\begin{eqnarray}\label{em}
K : \CE^{\mathrm{op}} \times \CP \lra \mathrm{Par}\mathrm{Set}
\end{eqnarray}
defined on objects by 
$$K(A,X) = \CA(A,X) \ .$$
It is defined on morphisms by
\begin{eqnarray}\label{strepi}
\begin{aligned}
\xymatrix{
K(A,X)=\CA(A,X) \ar[dd]_-{K(s,f)} \\
& \CA(A,U) \ar[lu]_-{\CA(A,f_0)} \ar[ld]^-{\CA(s,f_1)} \\
K(B,Y)=\CA(B,Y)}
\end{aligned}
\end{eqnarray}
where $s:B\lra A$ is in $\CE$ and $[f] = [f_0,U,f_1]:X\lra Y$ 
is a partial map in $\CA$. To see that composition is preserved, take $t:C\lra B$
in $\CE$ and a partial map $[g]=[g_0,V,g_1]:Y\lra Z$,
and form the pullback $P$ of $f_1$ and $g_0$ as in the diamond in the diagram: 
\begin{equation}\label{pmcomp}
\begin{aligned}
\xymatrix{
& & P \ar[ld]^-{p_0} \ar@/_2pc/[lldd]_-{h_0} \ar[rd]_-{p_1} \ar@/^2pc/[rrdd]^-{h_1} \\
& U \ar[ld]^-{f_0} \ar[rd]_-{f_1} & & V \ar[rd]_-{g_1} \ar[ld]^-{g_0} \\
X & & Y & & Z \ .}
\end{aligned}
\end{equation}
Then we have
\begin{equation}\label{Mcomp}
\begin{aligned}
\xymatrix{
& & \CA(A,P) \ar[ld]^-{\CA(A,p_0)} \ar@/_2pc/[lldd]_-{\CA(A,h_0)} \ar[rd]_-{\CA(s,p_1)} \ar@/^2pc/[rrdd]^-{\CA(st,h_1)} \\
& \CA(A,U) \ar[ld]^-{\CA(A,f_0)} \ar[rd]_-{\CA(s,f_1)} & & \CA(B,V) \ar[rd]_-{\CA(t,g_1)} \ar[ld]^-{\CA(B,g_0)} \\
\CA(A,X) & & \CA(B,Y) & & \CA(C,Z)}
\end{aligned}
\end{equation}
where the diamond is a pullback obtained as a pasting of two pullbacks, 
one pullback is \eqref{strepi} for $s\in \CE$ and $p_0\in \CM$, 
and the other is the image under $\CA(B,-)$ of the diamond pullback in \eqref{pmcomp}. 
That is,
$$K(t,[g])\circ K(s,[f])=[\CA(A,h_0),\CA(A,P),\CA(st,h_1)]=K(st,[g]\circ[f]) \ .$$
So $K$ preserves composition. That $K$ preserves identity morphisms is clear.

We write 
$$\overline{K} : \CE^{\mathrm{op}}  \lra [ \CP,\mathrm{Par}\mathrm{Set}]$$ 
for the functor corresponding to \eqref{em} using cartesian closedness of $\mathrm{CAT}$. On objects, $\overline{K}(A) = K(A,-)$.

Note that, in $\mathrm{Par}\mathrm{Set}$, we can identify morphisms
$[f] = [f_0,U,f_1]:X\lra Y$ with their canonical unique representative $f$
of the form $(f_0,U,f_1)$ where $U\subseteq X$ 
(called the {\em domain of definition} of $f$) and $f_0$ is the inclusion.   

There is a well known equivalence of categories
\begin{eqnarray}\label{par=pted}
1+- : \mathrm{Par}\mathrm{Set} \simeq 1/\mathrm{Set}
\end{eqnarray}
where $1/\mathrm{Set}$ is the category of pointed sets 
and point-preserving functions. The equivalence $1+-$ takes a set
$X$ to the set $1+X$ which is $X$ with a disjointly added point $0$ and takes
a partial function $f : X \lra Y$ to the function $1+f : 1+X \lra 1+Y$ defined by 
\begin{equation*}
(1+f)(x) =
\begin{cases}
0 & \text{for }  x\notin U \\
f(x) & \text{for } x \in U \ ,
\end{cases}
\end{equation*}
where $U$ is the domain of definition of $f$.

The equivalence \eqref{par=pted} makes it clear that $\mathrm{Par}\mathrm{Set}$
is complete and cocomplete. 
For later purposes, we remind the reader that coproducts in $1/\mathrm{Set}$
are formed by taking the disjoint union and then identifying all the distinguished
points $0$. Moreover, $1/\mathrm{Set}$ is a monoidal category whose tensor
product is the smash product $(1+S)\wedge (1+X) = 1 + S\times X$;
(this is not the categorical product in $1/\mathrm{Set}$).

Since $[\CE, 1/\mathrm{Set}]$ is the free cocompletion of the free $1/\mathrm{Set}$-enriched category on $\CE^{\mathrm{op}}$ and
since $[\CP,1/\mathrm{Set}]$ is a cocomplete $1/\mathrm{Set}$-enriched category, the functor $\overline{K}$ and the equivalence \eqref{par=pted}
determine a colimit-preserving (pointed) functor
\begin{eqnarray}\label{hat}
\widehat{(-)} : [\CE, 1/\mathrm{Set}] \lra [\CP, 1/\mathrm{Set}] \ . 
\end{eqnarray}
We now look in more detail at \eqref{hat} and its right adjoint 
\begin{eqnarray}\label{tilde}
\widetilde{(-)} : [\CP, 1/\mathrm{Set}] \lra [\CE, 1/\mathrm{Set}] \ . 
\end{eqnarray}  

\section{The transform adjunction}\label{ta}

We continue with a category $\CA$ and prefactorization system as in Section~\ref{km}; 
however we now assume $(\CE,\CM)$ is a factorisation system \cite{FreydKelly} for 
which each $m\in \CM$ is a monomorphism. 
So now every morphism $f : A \lra X$ in $\CA$ factors as $f=me$ with
$e\in \CE$ and $m\in \CM$. It follows that this factorisation is unique up to
isomorphism. Moreover, the $\CM$-subobjects of each $X$ is partially ordered:
the order on $\CM$-subobjects is induced by existence of a morphism in $\CA/X$.
It also follows that $g\circ f \in \CE$ implies $g \in \CE$.

We abuse notation by writing $U\le X$ for an $\CM$-subobject of $X$. 
If necessary, we feel free to choose a representative of the isomorphism class and
denote it by $i_U : U \lra X$.
For each $f : X \lra Y$ and $U\le X$, define $fU \le Y$ by $f \circ i_U=i_{fU}\circ e$ where $e\in \CE$. 
For each $f : X \lra Y$ and $V\le Y$, define $f^{-1}V\le X$ by taking $i_{f^{-1}V}$ to
be the pullback of $i_V$ along $f$. 

Morphisms of $\CP$ will be written as $(U,f) : X \lra Y$
where $U\le X$ and $f:U\lra Y$. 

We shall define a functor
\begin{eqnarray}\label{hat2}
\widehat{(-)} : [\CE, 1/\mathrm{Set}] \lra [\CP, 1/\mathrm{Set}] 
\end{eqnarray}
and show that it agrees with \eqref{hat}.

For $F : \CE \lra 1/\mathrm{Set}$ and $X\in \CP$, define
\begin{eqnarray}\label{hatformula}
\widehat{F}X = \sum_{A\le X}{FA}
\end{eqnarray}
where the sum is the coproduct of pointed sets.
For each $(U,f) : X \lra Y$ in $\CP$, define
\begin{eqnarray}
\widehat{F}(U,f) : \sum_{A\le X}{FA} \lra \sum_{B\le Y}{FB} 
\end{eqnarray} 
by taking its composite $\widehat{F}(U,f) \circ \mathrm{in}_A$ with the injection
at $A\le X$ to be $0$ unless $A\le U$, in which case it is the
composite $\mathrm{in}_B \circ Fs$ where $f \circ i_U = i_B \circ s$ with $s\in \CE$. 

\begin{proposition} Under the assumptions on $\CA$, $\CE$ and $\CM$, the functors 
\eqref{hat} and \eqref{hat2} are isomorphic.
\end{proposition}
\begin{proof}
Since \eqref{hat2} preserves colimits, it suffices to prove the functors agree on
objects of the form $F_Z=1+ \CE(Z,-)$. 
Using factorisation, we have
$$\widehat{F_K}X =1+ \sum_{A\le X}{\CE(Z,A)} 
\cong 1 +\CA(Z,A) \cong 1 + (\overline{K}Z)X \ , 
$$
showing that \eqref{hat} and \eqref{hat2} agree on such objects. 
We leave the check on morphisms to the reader. 
\end{proof}

\begin{corollary}\label{coendforhat} 
For all $F:\CE \lra 1/\mathrm{Set}$, there is a natural isomorphism
$$\widehat{F}X\cong \int^{A\in \CE}{(1+K(A,X))\wedge FA} \ .$$
\end{corollary}

Suppose we have a functor $T : \CP \lra 1/\mathrm{Set}$.
Let us say $x\in TX$ {\em is supported} when, for all 
$m: U\lra X$ in $\CM$, if $(Tm^*)x \neq 0$ then $m$ is invertible. 
Define 
\begin{eqnarray}\label{fs}
\widetilde{T}X = \{x\in TX : x \text{ is supported} \} \ .
\end{eqnarray}
Notice that $0$ is always supported.
Also, it suffices to test supportedness using inclusions for $m$.

\begin{lemma} For all $s : X \lra Y$ in $\CE$, the function 
$Ts_* : TX \lra TY$ restricts to give a function $\widetilde{T}s : \widetilde{T}X \lra \widetilde{T}Y$.  
\end{lemma}
\begin{proof}
We need to see that if $x\in TX$ is supported then so is $(Ts_*)x\in TY$. 
Take any $n:V\lra Y$ in $\CM$ such that $(Tn^*)(Ts_*)x \neq 0$.
Form the pullback
\begin{equation}\label{pbmse}
\xymatrix{
U \ar[rr]^-{q} \ar[d]_-{m} && V \ar[d]^-{n} \\
X \ar[rr]_-{s} && Y}
\end{equation}
in $\CA$. Using $(Tn^*)(Ts_*) = (Tq_*)(Tm^*)$, we see that 
$(Tm^*)x \neq 0$. Since $x$ is supported, $m$ is invertible.
So $nq \in \CE$ and hence $n\in \CE$.
So $n$ is invertible. This shows that $(Ts_*)x \in \widetilde{T}Y$.   
\end{proof}

\begin{theorem}\label{mainadj}
The functor \eqref{hat2} has a right adjoint taking the functor $T : \CP \lra 1/\mathrm{Set}$ to $\widetilde{T}$.
The unit of the adjunction is invertible; in other words, \eqref{hat2} is fully faithful.
The counit of the adjunction is a monomorphism. 
\end{theorem}
\begin{proof}
We must construct a natural bijection
\begin{eqnarray}\label{mainadjhoms}
[\CP, 1/\mathrm{Set}](\widehat{F},T) \cong [\CE, 1/\mathrm{Set}](F,\widetilde{T}) \ .
\end{eqnarray}
Consider a natural transformation $\theta : \widehat{F} \Lra T$. 
Let $\theta_{X,A} : FA \lra TX$ denote the restriction of $\theta_X$ along
the $(A\le X)$-component of the injection. 
Observe that $\theta_{A,A} : FA \lra TA$ lands in $\widetilde{T}A$. 
For, suppose we have a $U\le A$ with 
$(Ti_U^*)\theta_{A,A}a\neq 0$
invertible for some $a\in FA$. Then, by naturality,
$\theta_{A,U}(\widehat{F}i_U^*)a\neq 0$.
Of course, then $(\widehat{F}i_U^*)a\neq 0$, and by definition of
$\widehat{F}$ on morphisms, $A\le U$. So $U=A$, as required.

Now we can define $\phi_A : FA \lra \widetilde{T}A$ by $\phi_Aa=\theta_{A,A}a$.
We shall show that these morphisms are the components of a 
natural transformation $\phi : F \Lra \widetilde{T}$. Take any
$s:A\lra B$ in $\CE$. Then
$$(\widetilde{T}s)\phi_Aa = (Ts_*)\theta_{A,A}a = \theta_B(\widehat{F}s_*)\mathrm{in}_Aa=\theta_{B,B}(Fs)a = \phi_B(Fs)a$$
using the naturality of $\theta$. 
So we have a natural function $\theta \mapsto \phi$ for \eqref{mainadjhoms}.  

For the inverse bijection, take a natural transformation $\phi : F \lra \widetilde{T}$
and define $\theta_X : \widehat{F}X \lra TX$ by
$$\theta_{X,A}a = (Ti_{A *})\phi_Aa \ .$$ 
We prove naturality of $\theta$ in two steps, first with respect to morphisms 
$i_U^*$ for $U\le X$ and then with respect to $f_*$ for $f : X\lra Y$
in $\CA$.
For the first, form the pullback span $(u,W,v) : A \lra V$ of the pair 
$i_A : A \lra X$ and $i_U :U \lra X$. Then
$$(Ti_U^*)\theta_{X,A}a = (Ti_U^*)(Ti_{A *})\phi_Aa = (Tv_*)(Tu^*)\phi_Aa = T(v_*u^*)\phi_Aa$$
which is $0$ unless $(Tu^*)\phi_Aa \neq 0$. However, $\phi_Aa$ is supported,
so $u$ is invertible. On the other hand, 
$\theta_U(\widehat{F}i_U^*)\mathrm{in}_Aa = 0$
unless $A\le V$ with representative $j:A\lra V$, say, in which case 
$$\theta_U(\widehat{F}i_U^*)\mathrm{in}_Aa = \theta_{U,A}a = (Tj_*)\phi_Aa = T(v_*u^*)\phi_Aa \ .$$ 

Now we turn to naturality of $\phi$ with respect to $f_*$. This requires factoring
$f i_A = i_B s$ with $s\in \CE$. Then
\begin{align*}
\theta_Y(\widehat{F}f_*)a & = \theta_Y(Fs)a  = (Ti_{B*})\phi_B(Fs)a = (Ti_{B*})(\widetilde{T}s)\phi_Aa \\  
& = T(i_{B*}s_*)\phi_Aa = T(f_*i_{A*})\phi_Aa  = (Tf_*)\theta_Xa \ .
\end{align*}

It is easy to see that $\theta \mapsto \phi$ and $\phi \mapsto \theta$ are mutually inverse.  

The component of the unit $\eta_F : F \Lra \widetilde{\widehat{F}}$ at $A$ is
just the injection $\mathrm{in}_A : FA \lra \widehat{F}A$ seen as landing in
the supported elements; so it is injective.  
Take any supported $\mathrm{in}_Ux \in \widehat{F}A$. 
So $x\in FU$. If $x=0$ then it is in the image of $\eta_FA$.
If $x\neq 0$ then $(\widehat{F}i_U^*)x = x \neq 0$, and, by supportedness,
$U = A$; so $\mathrm{in}_Ux = \mathrm{in}_Ax$ is in the image of $\eta_FA$. 
So $\eta_F$ is invertible. 

The component of the counit $\varepsilon_T : \widehat{\widetilde{T}} \lra T$
at $X$ takes $a\in \widetilde{T}A$ to $(Ti_{A*})a \in TX$ for $A\le X$.  
If $(Ti_{A*})a = 0$ then $(Ti_A^*)(Ti_{A*})a = 0$; but $i_A^*i_{A*} = 1_A$,
so $a=0$. If $(Ti_{A*})a = (Ti_{B*})b \neq 0$ with $a\in \widetilde{T}A$,
$b\in \widetilde{T}B$, then $T(i_B^*i_{A*})a = T(i_B^*i_{B*})b = b$.
Form the pullback $(u,P,v)$ of $i_A$ and $i_B$. Then we have
$T(u_*v^*)a = T(i_B^*i_{A*})a = b \neq 0$. So $T(v^*)a \neq 0$.
Since $a$ is supported, $v$ is invertible. It follows that $A\le B$.
Symmetrically, $B\le A$. So $A = B$ and $a = b$.
This proves injectivity of $\varepsilon_TX$. So $\varepsilon_T$ is
a monomorphism.       
\end{proof}

\begin{remark}
The essential image of \eqref{hat2} consists of those $T$ with the property that
each $x\in TX$ is of the form $x=T(i_{U*})u$ for some $U\le X$ and
supported $u\in TU$.
\end{remark}

\begin{corollary}\label{endfortilde} 
For all $T:\CP \lra 1/\mathrm{Set}$, there is a natural isomorphism
$$\widetilde{T}A\cong \int_{X\in \CP}{[K(A,X), TX]} \ .$$
\end{corollary}

There is no problem replacing $1/\mathrm{Set}$ by 
the category $\mathrm{Mod}^R$. The same formula \eqref{hatformula}
provides a functor
\begin{eqnarray}\label{hatmod}
\widehat{(-)} : [\CE, \mathrm{Mod}^R] \lra [\CP, \mathrm{Mod}^R] \ .
\end{eqnarray}
Similarly, the formula \eqref{fs} gives the value at $T$ of a right adjoint 
\begin{eqnarray}\label{tildemod}
\widetilde{(-)} : [\CP, \mathrm{Mod}^R] \lra [\CE, \mathrm{Mod}^R] \ .
\end{eqnarray}
to \eqref{hatmod}. Notice that $\widetilde{T}X$ is a submodule of $TX$: supported
elements are closed under addition and scalar multiple
(if $(Tm^*)(x+y)= (Tm^*)x+(Tm^*)y \neq 0$ then either $(Tm^*)x \neq 0$ or $(Tm^*)y \neq 0$; and, if $(Tm^*)(rx)= r(Tm^*)x \neq 0$ then $(Tm^*)x \neq 0$).   

The proof of Theorem~\ref{mainadj} also proves:
\begin{theorem}\label{modadj}
The functor \eqref{hatmod} has the right adjoint \eqref{tildemod}.
The unit of the adjunction is invertible; in other words, \eqref{hatmod} is fully faithful.
The counit of the adjunction is a monomorphism. 
\end{theorem}


\section{Two comonads}

We continue with $\CA$, $\CE$, $\CM$ and $\CP$ as before.
  
Let $J:\CE \lra \CA$ be the inclusion functor and let $I = (-)_* :\CA \lra \CP$.
Both functors are the identity on objects.

Let $\CX$ be any category admitting coproducts and products indexed by
sets of $\CM$-subobjects of any given object of $\CA$.

\begin{lemma}\label{leftKan}
Each functor $F:\CE \lra \CX$ has a pointwise left Kan extension
\begin{equation}
\begin{aligned}
\xymatrix{
\CE \ar[rd]_{F}^(0.5){\phantom{a}}="1" \ar[rr]^{J}  && \CA \ar[ld]^{\mathrm{Lan}_JF}_(0.5){\phantom{a}}="2" \ar@{=>}"1";"2"^-{\kappa}
\\
& \CX 
}
\end{aligned}
\end{equation} 
along $J:\CE \lra \CA$ defined on objects by:
$$(\mathrm{Lan}_JF)X = \sum_{U\le X}FU \ .$$
For morphisms $f :X\lra Y$ in $\CA$, the following square commutes.
\begin{equation}
\begin{aligned}
\xymatrix{
FU \ar[rr]^-{\mathrm{in}_U} \ar[d]_-{Fe} && (\mathrm{Lan}_JF)X \ar[d]^-{(\mathrm{Lan}_JF)f} \\
FfU \ar[rr]_-{\mathrm{in}_{fU}} && (\mathrm{Lan}_JF)Y}
\end{aligned}
\end{equation}
\end{lemma}
\begin{proof}
The inclusion $U\le X \mapsto (U, i_U : JU\lra X)$ 
of the discrete category on the set $\{U : U\le X\}$ into the comma category
$J/X$ has a left adjoint, taking $(V,f:JV\lra X)$ to $fV\le X$, and so is final. 
It follows that the colimit of 
$J/X \stackrel{\mathrm{dom}}\lra \CE \stackrel{F}\lra \CX$
can be calculated by restricting along the inclusion and so is the coproduct displayed.
\end{proof}

\begin{lemma}
Each functor $T:\CA \lra \CX$ has a pointwise right Kan extension
\begin{equation}
\begin{aligned}
\xymatrix{
\CA \ar[rd]_{T}^(0.5){\phantom{a}}="1" \ar[rr]^{I}  && \CP \ar[ld]^{\mathrm{Ran}_IT}_(0.5){\phantom{a}}="2" \ar@{<=}"1";"2"^-{\rho}
\\
& \CX 
}
\end{aligned}
\end{equation}  
along $I:\CA \lra \CP$ defined on objects by:
$$(\mathrm{Ran}_IT)X = \prod_{U\le X}TU \ .$$
For morphisms $(W,h) :X\lra Y$ in $\CP$, the square
\begin{equation}
\begin{aligned}
\xymatrix{
(\mathrm{Ran}_IT)X \ar[rr]^-{\mathrm{pr}_{h^{-1}V}} \ar[d]_-{(\mathrm{Ran}_IT)(W,h)} && Th^{-1}V \ar[d]^-{T\bar{h}} \\
(\mathrm{Ran}_IT)Y \ar[rr]_-{\mathrm{pr}_V} && TV 
}
\end{aligned}
\end{equation}
commutes, where $\bar{h}$ is the pullback of $h$ along $i_V$.
\end{lemma}
\begin{proof}
The functor $U\le X \mapsto (i_U^* : X\lra IU, U)$ from the discrete category 
on the set $\{U : U\le X\}$ into the comma category
$X/I$ has a right adjoint, taking $((W,h): X\lra IA,A)$ to $W\le X$, and so is initial. 
It follows that the limit of 
$X/I \stackrel{\mathrm{cod}}\lra \CA \stackrel{T}\lra \CX$
can be calculated as the product displayed.
\end{proof}

This gives the two adjunctions
\begin{equation}\label{twoadj} 
\xymatrix @R-3mm {
[\CE,\CX] \ar@<1.5ex>[rr]^{\mathrm{Lan}_J} \ar@{}[rr]|-{\bot} && [\CA,\CX]
\ar@<1.5ex>[rr]^{\mathrm{Ran}_I}  \ar@{}[rr]|-{\top} \ar@<1.5ex>[ll]^{[J,1]} && [\CP,\CX]
\ar@<1.5ex>[ll]^{[I,1]}  \ .
}
\end{equation}
Each adjunction generates a comonad on $[\CA,\CX]$:
\begin{equation}\label{twocmnd}
G = \mathrm{Lan}_J \circ [J,1] \ \text{ and } \  H = [I,1] \circ \mathrm{Ran}_I \ .
\end{equation}
To be more explicit, the functors $GF$ and $HF$ are defined on objects by 
$$(GF)A=\sum_{U\le A}{FU} \ \text{ and } \ (HF)A=\prod_{U\le A}{FU} \ .$$ 
On morphisms $f:A\lra B$ in $\CA$, they are defined by commutativity in the squares
\begin{equation}
\begin{aligned}
\xymatrix{
FU \ar[rr]^-{Fe} \ar[d]_-{\mathrm{in}_U} && FfU \ar[d]^-{\mathrm{in}_{fU}} \\
(GF)A \ar[rr]_-{(GF)f} && (GF)B}
\qquad
\xymatrix{
(HF)A \ar[rr]^-{(HF)f} \ar[d]_-{\mathrm{pr}_{f^{-1}V}} && (HF)B \ar[d]^-{\mathrm{pr}_V} \\
Ff^{-1}V \ar[rr]_-{F\bar{f}} && FV \ \ .}
\end{aligned}
\end{equation}
The counits $\varepsilon_F : GF \lra F$ and $\varepsilon_F : HF \lra F$ have
components respectively defined by $\varepsilon_FA \circ \mathrm{in}_U = F\mathrm{in}_U$
and $\varepsilon_FA = \mathrm{pr}_A$. 
The comultiplications $\delta_F : GF\lra G^2F$ and $\delta_F : HF\lra H^2F$ 
have components respectively defined as follows.
\begin{equation*}
\xymatrix{
FU \ar[rd]_{\mathrm{in}_{U\le U}}\ar[rr]^{\mathrm{in}_U}   && \sum_{U\le A}{FU} \ar[ld]^{\delta_FA} \\
& \sum_{V\le U\le A}{FV}  &
}
\qquad
\xymatrix{
\prod_{U\le A}{FU} \ar[rd]_{\delta_FA}\ar[rr]^{\mathrm{pr}_{V}}   && FV  \\
& \prod_{V\le U\le A}{FV}\ar[ru]_{\mathrm{pr}_{V\le U}}  &
}
\end{equation*}


We shall construct a comonad morphism $\Theta : G \Lra H$ when $\CX$ is a pointed
category.

Take $F : \CA \lra \CX$ and $A\in \CA$. We define 
\begin{equation}\label{Thetacmpnt}
\Theta_FA : \sum_{U\le A}FU \lra \prod_{V\le A}FV
\end{equation}
by taking the composite
$$ FU\stackrel{\mathrm{in}_U}\lra \sum_{U\le A}{FU} \stackrel{\Theta_FA}\lra \prod_{V\le A}{FV} \stackrel{\mathrm{pr}_V}\lra FV$$
to be zero unless $U\le V$, in which case the composite is $F(U\le V) : FU \lra FV$.

\begin{proposition}
For the two comonads $G$ and $H$ of \eqref{twocmnd} on $[\CA,\CX]$, a comonad morphism $\Theta : G \Lra H$ is defined by the components \eqref{Thetacmpnt}.
\end{proposition}
\begin{proof}
Naturality in $A$ is proved by contemplating the following diagram.
\begin{equation*}
\xymatrix{
FU \ar[d]_{Fe}\ar[r]^{\mathrm{in}_U \ }  & \sum_{U\le A}{FU}\ar[d]^{(GF)f} \ar[r]^{\Theta_FA} & \prod_{U_1\le A}{FU_1} \ar[d]^{(HF)f} \ar[r]^{ \ \mathrm{pr}_{f^{-1}V}} 
& Ff^{-1}V \ar[d]^{F\bar{f}} \\
FfU\ar[r]_{\mathrm{in}_{fU} \ }  & \sum_{V_1\le B}{FV_1}\ar[r]_{\Theta_FB} & \prod_{V\le B}{FV}\ar[r]_{ \ \mathrm{pr}_V} 
& FV }
\end{equation*}
The top composite is zero unless $U\le f^{-1}V$ (say with inclusion $i$). 
The bottom composite is zero unless $fU\le V$ (say with inclusion $j$). 
The conditions are the same.
When this condition holds, the diagram commutes since $\bar{f} \circ i = j \circ e$.  

Naturality in $F$ is obvious. 

Preservation of the counits is proved by the calculation
$$\varepsilon_FA\circ \Theta_FA \circ \mathrm{in}_U = \mathrm{pr}_A\circ \Theta_FA \circ \mathrm{in}_U
= F\mathrm{in}_U = \varepsilon_FA\circ \mathrm{in}_U \ .$$ 

Preservation of the comultiplications is proved by observing that 
$$\mathrm{pr}_{V\le U} \circ \delta_FA \circ \Theta_FA \circ \mathrm{in}_W
= \mathrm{pr}_{V} \circ \Theta_FA \circ \mathrm{in}_W$$
while
$$\mathrm{pr}_{V\le U} \circ \Theta^2_FA \delta_FA  \circ \mathrm{in}_W
= \mathrm{pr}_{V\le U} \circ \Theta^2_FA \circ \mathrm{in}_{W\le W} \ .$$
Both of these right-hand sides are zero unless $W\le V$, in which case they
are both equal to $F(W\le V)$. 
\end{proof}

\begin{proposition}
The functor $[I,1] : [\CP,\CX]\lra [\CA,\CX]$ is comonadic.
\end{proposition}
\begin{proof}
The functor $[I,1]$ is conservative since $I$ is bijective on objects.
It also preserves all limits, including all equalizers.
The result follows (for example) by the Beck comonadicity theorem \cite{ML1963}.
\end{proof}

It follows that $\Theta :G\Lra H$
induces a functor $\bar{\Theta} : [\CA,\CX]^G \lra [\CP,\CX]$ over $[\CA,\CX]$, 
where $[\CA,\CX]^G$ is the category of Eilenberg-Moore $G$-coalgebras.  
The composite of $\bar{\Theta}$ with the comparison functor 
$[\CE,\CX] \lra [\CA,\CX]^G$ is isomorphic to  
\begin{eqnarray}\label{hatX}
\widehat{(-)} : [\CE, \CX] \lra [\CP, \CX]
\end{eqnarray}
as defined in \eqref{hatformula}.

\section{The Morita equivalence}

With a finite well-poweredness assumption on the factorization system of $\CA$, 
we shall show that the adjunction of Theorem~\ref{modadj} is an equivalence.
In other words, we shall show that the free additive categories on
$\CE$ and $\CP$ are Morita equivalent (= additively Cauchy equivalent).

Our goal in this section is to prove our main theorem.

\begin{theorem}\label{mainthm}
Let $(\CE,\CM)$ be a factorisation system on a category $\CA$.
Assume that $\CM$ is contained in the class of monomorphisms
and that all pairs $(m :U\rightarrow X,f:A\rightarrow X)$ of morphisms with $m\in \CM$
have a pullback in $\CA$. 
Assume that the $\CM$-subobjects of each object form a finite set.
Regard $\CE$ as a subcategory of $\CA$ with the same objects
and put $\CP = \mathrm{Par}_{\CM}\CA$. 
Let $\CX$ be any finitely complete additive category.
Then the functor \eqref{hatX} is an equivalence of categories
$$[\CE, \CX] \simeq [\CP, \CX] \ . $$
\end{theorem}

We first prove that the comonads $G$ and $H$ of \eqref{twocmnd} are isomorphic.

\begin{proposition}\label{invTheta}
Under the assumptions of Theorem~\ref{mainthm}, the comonad morphism
$\Theta : G\Lra H$, with components \eqref{Thetacmpnt}, is invertible.
\end{proposition}
\begin{proof}
Since $\CX$ is additive, every finite product is a coproduct.
The ordered set of $\CM$-subobjects of $A \in \CA$ is assumed finite
and so has a linear refinement. 
So we can list all the subobjects as $0=U_0, U_1, \dots , U_n = A$ 
such that $0\le i \le j \le n$ implies $U_i \le U_j$.
Then the morphisms \eqref{Thetacmpnt} are represented by an
upper-triangular matrix with identity morphisms in the main diagonal.
Since $\CX$ is additive (including existence of additive inverses), such
matrices are invertible.
\end{proof}

\begin{proof}[Proof of Theorem~\ref{mainthm}]
By Proposition~\ref{invTheta}, the functor \eqref{hatX} becomes the
comparison from $[\CE,\CX]$ into the category of $G$-coalgebras. 
So the Theorem is now equivalent to comonadicity of the functor 
$\mathrm{Lan}_J : [\CE,\CX]\lra [\CA,\CX]$.
We already know that \eqref{hatX} is fully faithful.
So it remains to prove that $\mathrm{Lan}_J$ preserves certain equalizers.
In fact, it preserves all finite limits. Since limits in $[\CA,\CX]$ are formed
pointwise, it suffices to see that each 
$$\mathrm{ev}_A\circ \mathrm{Lan}_J : [\CE,\CX]\lra \CX$$
preserves finite limits where $\mathrm{ev}_A : [\CA,\CX]\lra \CX$ is evaluation
at $A\in \CA$. By Lemma~\ref{leftKan}, 
$$\mathrm{ev}_A\circ \mathrm{Lan}_J \cong \bigoplus_{U\le A}\mathrm{ev}_U \ .$$
Each evaluation $\mathrm{ev}_U$ preserves limits so their direct sum does.    
\end{proof}

\section{Some examples}\label{exx}

A {\em (set) species} in the sense of Joyal \cite{Species} is a functor 
$F : \mathfrak{S} \lra \mathrm{Set}$
where $\mathfrak{S}$ is the category of sets and bijective functions. 
A {\em pointed-set species} is a functor 
$F : \mathfrak{S} \lra 1/\mathrm{Set}$.
An {\em $R$-module species} is a functor 
$F : \mathfrak{S} \lra \mathrm{Mod}^R$; the case where $R$ is
a field is the basic situation of \cite{AnalFunct}.

Following \cite{CEF2012}, we write $\mathrm{FI}$ for the category of finite sets
and injective functions. We then see that 
$\CM = \mathrm{FI}$ and $\CE = \mathfrak{S}$,
while $\CP = \mathrm{FI}\sharp$, the category of injective partial functions.  

\begin{corollary} \cite{CEF2012} The functor
$$\widehat{(-)} : [\mathfrak{S}, \mathrm{Mod}^R] \lra [\mathrm{FI}\sharp, \mathrm{Mod}^R]$$
is an equivalence of categories.
\end{corollary}

Another example relevant to \cite{repfgl} is the category $\CA = \mathrm{FIVect}_{\mathbb{F}_q}$
of finite vector spaces over the field $\mathbb{F}_q$ of cardinality 
$q$ (a prime power) and injective linear functions.
Let $\CE = \mathfrak{Gl}_q$ be the category of finite
$\mathbb{F}_q$-vector spaces and linear isomorphisms. 
Then $\CP = \mathrm{FI}\sharp_q$ is the
category of finite $\mathbb{F}_q$-vector spaces and injective 
partial linear functions.

\begin{corollary} The functor
$$\widehat{(-)} : [\mathfrak{Gl}_q, \mathrm{Mod}^R] \lra [\mathrm{FI}\sharp_q, \mathrm{Mod}^R]$$
is an equivalence of categories.
\end{corollary}   

Here are a few examples of categories $\CA$ to which 
Theorem~\ref{mainthm} applies with $\CE$ 
the epimorphisms and $\CM$ the monomorphisms:
\begin{enumerate}[(a)]
\item the category of finite abelian groups and group morphisms;
\item the category of finite abelian $p$-groups and group morphisms;
\item the category of finite sets and all functions.
\end{enumerate}
Theorem~\ref{mainthm} also applies to $\CM$ in place of $\CA$
in these examples. Then $\CE$ is replaced by the groupoid of
invertible morphisms in $\CA$. In case of example (a),
the paper \cite{HillarDarren2007} describes the groupoid being represented in $\CX$. 

Consider the ``algebraic'' simplicial category $\DelCat$ 
whose objects are all the natural numbers and 
whose morphisms $\xi : m\lra n$ are order-preserving functions 
$$\xi : \{0, 1, \dots , m-1\}\lra \{0, 1, \dots , n-1\} \ .$$
Put $\CA = \DelCat^{\mathrm{op}}$.
Take $\CM$ in $\CA$ to consist of the surjections in $\DelCat$.
Pushouts of surjections along arbitrary morphisms exist in 
$\DelCat$. 
Then $\CE = \DelCat_{\mathrm{inj}}^{\mathrm{op}}$
and $\CP$ is the opposite of the category whose morphisms
$m\lra n$ are cospans $$m\stackrel{\xi}\lra r \stackrel{\sigma}\lla n$$
in $\DelCat$ with $\sigma$ surjective. 
We could also take the ``topological'' simplicial category 
$\DelCat_{\mathrm{t}}$ (omit the object $0$) to obtain
a reinterpretation of the preoperads in $\CX$ in the sense of 
\cite{Berger1996}.

Here is a rather trivial example involving $\DelCat_{\mathrm{t}}$.
Take $\CA$ to be the category of non-empty ordinals and morphisms
the order-preserving functions which preserve first element.
Let $\CM$ be the class of morphisms which are inclusions of
initial segments. This is part of a factorization system where
$\CE = \DelCat_{\bot \neq \top}$ is the category of ordinals with 
distinct first and last element 
and morphisms the order-preserving functions which preserve first 
and last element. Sometimes $\CE$ is called the category of {\em intervals}; there is a duality isomorphism
$$\CE \cong \DelCat_{\mathrm{t}}^{\mathrm{op}} \ .$$ 
In this case, not only do we have the equivalence
$$[\CE, \CX] \simeq [\CP, \CX] $$
of Theorem~\ref{mainthm}, we actually have an isomorphism
$$\CP \cong \CE \ .$$



Take $\CA$ to be a (partially) ordered set with finite infima and the descending chain condition. 
Then every morphism is a monomorphism
and the strong epimorphisms are equalities. So $\CE$ is
the discrete category $\mathrm{ob}\CA$ on the set of
elements of the ordered set. The reader may like to investigate
the case where $\CA$ is the set of strictly positive integers
ordered by division. 

\section{Appendix: Remarks on the Dold-Puppe-Kan equivalence}\label{DPK}

Let $\mathbb{D}$ be the additive category whose objects are the natural numbers,
whose homs are
\begin{equation*}
\mathbb{D}(m,n) =
\begin{cases}
\mathbb{Z} & \text{for } n=m,m+1 \ , \\
0 & \text{otherwise } \ ,
\end{cases}
\end{equation*}
and whose composition
$$\mathbb{D}(n,p)\otimes \mathbb{D}(m,n)\lra \mathbb{D}(m,p)$$
is zero unless all three homs are $\mathbb{Z}$ in which case it
is the canonical isomorphism $\mathbb{Z}\otimes \mathbb{Z}\cong \mathbb{Z}$.
We write $\partial_n : n\lra n+1$ and $1_n : n\lra n$ for $1\in \mathbb{Z}$.

For any additive category $\CV$, the additive functor category 
$[\mathbb{D}^{\mathrm{op}},\CV]$
is the category of differential non-negatively graded $\CV$-objects
(or $\CV$-chain complexes).

Let $\DelCat_{\mathrm{inj}}$ denote the monoidal subcategory of the algebraic simplicial category $\DelCat$ whose objects are all the natural numbers and whose morphisms $m\lra n$ are order-preserving injective functions 
$$\mu : \{0, 1, \dots , m-1\}\lra \{0, 1, \dots , n-1\} \ .$$
This is the PRO for coaugmented objects: that is, objects equipped with
a morphism from the tensor unit. Let $K:\DelCat_{\mathrm{inj}}\lra \DelCat$ 
denote the inclusion.

There is a functor $J: \DelCat_{\mathrm{inj}} \lra \mathbb{D}$ which
is the identity on objects and takes the injection $\partial_i : n-1\lra n$
to $0$ for $i<n$ and to $\partial_n$ for $i=n$.

At the heart of the Dold-Puppe-Kan Theorem is the span of functors
$$\mathbb{D}\stackrel{J}\lla \DelCat_{\mathrm{inj}}\stackrel{K}\lra \DelCat \ .$$ 

For any simplicial object $X: \DelCat\lra \CV$, consider the
chain complex $N(X)$ defined as the right Kan extension of 
$X K^{\mathrm{op}}$ along $J^{\mathrm{op}}$, 
which exists under a mild completeness condition on $\CV$.
\begin{equation}
\begin{aligned}
\xymatrix{
\DelCat_{\mathrm{inj}}^{\mathrm{op}} \ar[d]_{K^{\mathrm{op}}}^(0.5){\phantom{aaa}}="1" \ar[rr]^{J^{\mathrm{op}}}  && \mathbb{D}^{\mathrm{op}}  \ar[d]^{N(X)}_(0.5){\phantom{aaa}}="2" \ar@{<=}"1";"2"^-{\rho}
\\
\DelCat^{\mathrm{op}}  \ar[rr]_-{X} && \CV 
}
\end{aligned}
\end{equation} 
We have the weighted limit formula
$$N(X)_n=\mathrm{lim}(\mathbb{D}(J-,n),XK^{\mathrm{op}}) \ .$$
The weight functor $\mathbb{D}(J-,n): \DelCat_{\mathrm{inj}}^{\mathrm{op}} \lra \mathrm{Ab}$
into the category $\mathrm{Ab}$ of abelian groups takes all objects to $0$
except that $n-1$ and $n$ go to $\mathbb{Z}$, and takes the morphisms
$\partial_i : n-1\lra n$ to $0$ for $0\le i <n$ and to the identity abelian group morphism $\mathbb{Z}\lra \mathbb{Z}$ for $i=n$. 
The natural transformations $\mathbb{D}(J-,n) \Lra SK^{\mathrm{op}}$ 
for a simplicial abelian group $S$ are clearly in bijection with elements
$x\in S_n$ such that $d_i(x) = 0$ for $0\le i < n$. 
It follows that $N(X)_n$ is the intersection of the kernels of the 
$d_i: X_n\lra X_{n-1}$ for $0\le i < n$.
Also, $d_n : N(X)_n\lra N(X)_{n-1}$ is induced by $d_n : X_n\lra X_{n-1}$,
making $N(X)$ a chain complex in $\CV$.    

We have the composite adjunction
\[ \xymatrix @R-3mm {
[\mathbb{D}^{\mathrm{op}},\CV] \ar@<1.5ex>[rr]^{\mathrm{Res}_{J^{\mathrm{op}}}} \ar@{}[rr]|-{\perp} && [\DelCat_{\mathrm{inj}}^{\mathrm{op}},\CV] \ar@<1.5ex>[ll]^{\mathrm{Ran}_{J^{\mathrm{op}}}} \ar@<1.5ex>[rr]^{\mathrm{Lan}_{K^{\mathrm{op}}}} \ar@{}[rr]|-{\perp} && [\DelCat^{\mathrm{op}},\CV] \ar@<1.5ex>[ll]^{\mathrm{Res}_{K^{\mathrm{op}}}} 
}\]
where $\mathrm{Res}_{H}$ denotes restriction along $H$ and its adjoints are
left and right Kan extension along $H$. In our cases of $H$, the Kan extensions
exist when $\CV$ is finitely complete (and so has a bit of cocompleteness since
it has finite direct sums).

Above we have identified the composite right-adjoint functor $\mathrm{Ran}_{J^{\mathrm{op}}} \mathrm{Res}_{K^{\mathrm{op}}}$ as the normalization functor $N$
from simplicial objects to chain complexes. 
It is this functor that is proved an equivalence by Dold-Puppe-Kan~\cite{Dold, DoldPuppe, Kan}. 

Since $K$ is strong monoidal, it follows from \cite{DaySt1995} that
$\mathrm{Lan}_{K^{\mathrm{op}}}$ is strong monoidal for the convolution
monoidal structures on $[\DelCat^{\mathrm{op}},\CV]$ and  $[\DelCat_{\mathrm{inj}}^{\mathrm{op}},\CV]$.   

Although I have not completed the calculation, it seems that 
$\mathrm{Res}_{J^{\mathrm{op}}}$ is also strong monoidal for convolution
monoidal structures on $[\DelCat_{\mathrm{inj}}^{\mathrm{op}},\CV]$
and the usual tensor product on $[\mathbb{D}^{\mathrm{op}},\CV] $. 
Since $\mathrm{Res}_{J^{\mathrm{op}}}$ is colimit preserving, we
only need check the result on representables.   

This would imply that $eN : [\DelCat^{\mathrm{op}},\CV]\lra [\mathbb{D}^{\mathrm{op}},\CV]$ is a monoidal equivalence. 

The monoidal structure on $\DelCat$ is not symmetric and so neither is the
monoidal structure on $[\DelCat^{\mathrm{op}},\mathrm{Set}]$.
As often happens, when we move to the additive context, a symmetry can be
found: just transport the usual symmetry on chain complexes across the equivalence $N$. 

There is something tantalizingly similar about the Dold-Puppe-Kan equivalence and the equivalence of Theorem~\ref{mainthm}.
We mentioned towards the end of Section~\ref{exx} that there is
an isomorphism
$$\DelCat_{\mathrm{t}}^{\mathrm{op}} \cong \DelCat_{\bot \neq \top} \ .$$  
From this viewpoint, the Dold-Puppe-Kan equivalence 
$$N : [\DelCat_{\bot \neq \top},\CX] \lra [\mathbb{D}^{\mathrm{op}},\CX]$$
is defined by
$$N(T)n = \bigcap_{0\le i < n}{\mathrm{ker}(T(\sigma_i) : T(n+1)\lra Tn)} \ .$$
The formula we have been working with for the right adjoint
$$\widetilde{(-)} : [\CP, \CX] \lra [\CE, \CX]$$
in the equivalence of Theorem~\ref{mainthm} can be expressed as
$$\widetilde{T}A = \bigcap_{\text{non-invertible }m\in \CM}{\mathrm{ker}( T(m^*) : TA\lra TU)} \ .$$  

Surely these results are special cases of something more general.  

\begin{thebibliography}{00}

\bibitem{Berger1996} Clemens Berger,\textit{Op\'erades cellulaires et espaces de lacets it\'er\'es}, Annales de L'Institut Fourier \textbf{46(4)} (1996) 1125--1157.\label{Berger1996} 

\bibitem{CEF2012} Thomas Church, Jordan S. Ellenberg and Benson Farb \textit{FI-modules: a new approach to stability for $S_n$-representations}, (\url{arXiv:1204.4533v2}, 28 Jun 2012).\label{CEF2012}

\bibitem{DaySt1995} Brian J. Day and Ross Street, \textit{Kan extensions along promonoidal functors}, Theory and Applications of Categories \textbf{1(4)} (1995) 72--77. \label{DaySt1995} 

\bibitem{DaySt1997} Brian J. Day and Ross Street, \textit{Monoidal bicategories and Hopf algebroids}, Advances in Math. \textbf{129} (1997) 99--157. \label{DaySt1997}

\bibitem{Dold} Albrecht Dold, \textit{Homology of symmetric products and other functors of complexes}, Annals of Math. \textbf{68} (1958) 54--80.\label{Dold}
 
\bibitem{DoldPuppe} Albrecht Dold and Dieter Puppe, \textit{Homologie nicht-additiver Funktoren. Anwendungen}, Ann. Inst. Fourier Grenoble \textbf{11} (1961) 201--312.\label{DoldPuppe}

\bibitem{FreydKelly} Peter J. Freyd and G. Max Kelly, \textit{Categories of continuous functors I}, J. Pure and Applied Algebra \textbf{2} (1972) 169--191; Erratum Ibid. \textbf{4} (1974) 121.\label{FreydKelly}

\bibitem{HillarDarren2007} Christopher J. Hillar and Darren L. Rhea, \textit{Automorphisms of finite abelian groups}, American Mathematical Monthly \textbf{114(10)} (2007) 917--923; \url{arXiv:math/0605185}.\label{HillarDarren2007}

\bibitem{Species} Andr\'e Joyal, \textit{Une theory combinatoire des series formelles}, Advances in Mathematics \textbf{42} (1981) 1--82.\label{Species}

\bibitem{AnalFunct} Andr\'e Joyal, \textit{Foncteurs analytiques et esp\`eces de structures}, Lecture Notes in Mathematics \textbf{1234} (Springer 1986) 126--159.\label{AnalFunct} 

\bibitem{TYBO} Andr\'e Joyal and Ross Street, \textit{Tortile Yang-Baxter operators in tensor categories}, J. Pure Appl. Algebra \textbf{71} (1991) 43--51.\label{TYBO}

\bibitem{ITD&QG} Andr\'e Joyal and Ross Street, \textit{An introduction to Tannaka duality and quantum groups}; Part II of Category Theory, Proceedings, Como 1990 (Editors A. Carboni, M.C. Pedicchio and G. Rosolini) Lecture Notes in Math. \textbf{1488} (Springer-Verlag Berlin, Heidelberg 1991) 411--492.\label{ITD&QG}

\bibitem{repfgl} Andr\'e Joyal and Ross Street, \textit{The category of representations of the general linear groups over a finite field}, J. Algebra \textbf{176 (3)} (1995) 908--946.\label{repfgl}

\bibitem{Kan} Daniel Kan, \textit{Functors involving c.s.s complexes}, Transactions of the American Mathematical Society \textbf{87} (1958) 330--346.\label{Kan}

\bibitem{Lindner1976} Harald Lindner, \textit{A remark on Mackey functors}, Manuscripta Mathematica \textbf{18} (1976) 273--278.\label{Lindner1976}

\bibitem{ML1963} Saunders Mac Lane, \textit{Natural associativity and commutativity}, Rice University Studies \textbf{49} (1963) 28--46.\label{ML1963}

\bibitem{NHA} Bodo Pareigis, \textit{A non-commutative non-cocommutative Hopf algebra in nature}, Journal of Algebra \textbf{70} (1981) 356--374.\label{NHA}

\bibitem{QG} Ross Street, \textit{Quantum Groups: A Path to Current Algebra}, Australian Mathematical Society Lecture Series \textbf{19} (Cambridge University Press, 2007).\label{QG} 

\bibitem{CBA} D. Tambara, \textit{The coendomorphism bialgebra of an algebra}, Journal of The Faculty of Science, The University of Tokyo, Section IA, Mathematics, \textbf{37} (1990) 425--456.\label{CBA}

\bibitem{Ulbrich} K.-H. Ulbrich, \textit{On Hopf algebras and rigid monoidal categories}, Israel Journal of 
Math. \textbf{72} (1990) 252--256.\label{Ulbrich}

\end{thebibliography}

\end{document}